% GENERATED FILE. Edit canonical lessons, then run npm run generate:books.
% canonical-source-hash: fnv1a64:d6ce64a8c44b8e0e
% canonical-lessons: FR-C310-adolescence, FR-C310-adolescent, FR-C310-jeunesse, FR-C310-rentree, FR-C310-majeur

\chapter{La jeunesse}
\label{ch:fr-la-jeunesse}
\hlchaptermodality{310}

\begin{chapteropening}
\textbf{By the end of this chapter:} \emph{I can talk about growing up: adolescence, teenagers, youth, the start of the school year and coming of age.}

Complete the last lesson of chapter 310: I can talk about growing up: adolescence, teenagers, youth, the start of the school year and coming of age.
\end{chapteropening}

\section[l'adolescence]{l'adolescence — adolescence, the teenage years}

\label{lesson:FR-C310-adolescence}



\begin{quote}
Before the new one: say the French for the pushchair, then the French for childhood.
\end{quote}

\subsection*{You'll want to know: l'adolescence}
\textbf{l'adolescence} — \textquotedblleft{}adolescence, the teenage years\textquotedblright{}.

\begin{quote}
Son adolescence a été difficile.
\end{quote}

\begin{quote}
\emph{His teenage years were difficult.}
\end{quote}

Feminine: \emph{une adolescence heureuse}.

\begin{cousinweb}
\emph{adolescence} is a learned borrowing from Latin \emph{adulēscentia} \textquotedblleft{}youth\textquotedblright{} (EtymDB), from \emph{adolēscere} \textquotedblleft{}to grow up\textquotedblright{}, built on \emph{alere} \textquotedblleft{}to nourish\textquotedblright{} (Lewis \& Short).
\end{cousinweb}

\begin{grammarlens}[title={What you've built}]
One more word for telling, explaining and describing at B1.
\end{grammarlens}

\subsection*{Your turn}
\begin{itemize}
\raggedright
  \item \emph{Say it:} \emph{l'adolescence}
  \item \emph{Say it:} \emph{l'adolescence}, once more
  \item \emph{Say it:} a sentence of your own with \emph{l'adolescence}
  \item \emph{Recall:} say the French for the pushchair, then the French for childhood, then say \emph{l'adolescence} again
\end{itemize}

\begin{tcolorbox}[breakable,colback=teal!4,colframe=teal!35!black,title={Before you move on}]
What does l'adolescence mean? (\textquotedblleft{}Adolescence, the teenage years\textquotedblright{}.) Say it once more.
\end{tcolorbox}
\section[l'adolescent]{l'adolescent — the teenager, the adolescent}

\label{lesson:FR-C310-adolescent}



\begin{quote}
Before the new one: say the French for childhood, then the French for adolescence.
\end{quote}

\subsection*{You'll want to know: l'adolescent}
\textbf{l'adolescent} — \textquotedblleft{}the teenager, the adolescent\textquotedblright{}.

\begin{quote}
Les adolescents passent des heures sur leur téléphone.
\end{quote}

\begin{quote}
\emph{Teenagers spend hours on their phones.}
\end{quote}

Often \emph{un ado}, \emph{une ado} in speech.

\begin{cousinweb}
\emph{adolescent} is a learned borrowing from Latin \emph{adolēscēns} \textquotedblleft{}growing up\textquotedblright{} (EtymDB), the present participle of \emph{adolēscere}.
\end{cousinweb}

\begin{grammarlens}[title={What you've built}]
One more word for telling, explaining and describing at B1.
\end{grammarlens}

\subsection*{Your turn}
\begin{itemize}
\raggedright
  \item \emph{Say it:} \emph{l'adolescent}
  \item \emph{Say it:} \emph{l'adolescent}, once more
  \item \emph{Say it:} a sentence of your own with \emph{l'adolescent}
  \item \emph{Recall:} say the French for childhood, then the French for adolescence, then say \emph{l'adolescent} again
\end{itemize}

\begin{tcolorbox}[breakable,colback=teal!4,colframe=teal!35!black,title={Before you move on}]
What does l'adolescent mean? (\textquotedblleft{}The teenager, the adolescent\textquotedblright{}.) Say it once more.
\end{tcolorbox}
\section[la jeunesse]{la jeunesse — youth, young people}

\label{lesson:FR-C310-jeunesse}



\begin{quote}
Before the new one: say the French for adolescence, then the French for the teenager.
\end{quote}

\subsection*{You'll want to know: la jeunesse}
\textbf{la jeunesse} — \textquotedblleft{}youth, young people\textquotedblright{}.

\begin{quote}
Dans sa jeunesse, il jouait au rugby.
\end{quote}

\begin{quote}
\emph{In his youth, he played rugby.}
\end{quote}

Also young people as a group: \emph{la jeunesse d'aujourd'hui}.

\begin{cousinweb}
\emph{jeunesse} is built on \emph{jeune} \textquotedblleft{}young\textquotedblright{}, which is inherited from Latin \emph{iuvenis} \textquotedblleft{}young\textquotedblright{} (Lewis \& Short).
\end{cousinweb}

\begin{grammarlens}[title={What you've built}]
One more word for telling, explaining and describing at B1.
\end{grammarlens}

\subsection*{Your turn}
\begin{itemize}
\raggedright
  \item \emph{Say it:} \emph{la jeunesse}
  \item \emph{Say it:} \emph{la jeunesse}, once more
  \item \emph{Say it:} a sentence of your own with \emph{la jeunesse}
  \item \emph{Recall:} say the French for adolescence, then the French for the teenager, then say \emph{la jeunesse} again
\end{itemize}

\begin{tcolorbox}[breakable,colback=teal!4,colframe=teal!35!black,title={Before you move on}]
What does la jeunesse mean? (\textquotedblleft{}Youth, young people\textquotedblright{}.) Say it once more.
\end{tcolorbox}
\section[la rentrée]{la rentrée — the start of the school year; the return}

\label{lesson:FR-C310-rentree}



\begin{quote}
Before the new one: say the French for the teenager, then the French for youth.
\end{quote}

\subsection*{You'll want to know: la rentrée}
\textbf{la rentrée} — \textquotedblleft{}the start of the school year; the return\textquotedblright{}.

\begin{quote}
La rentrée a lieu en septembre.
\end{quote}

\begin{quote}
\emph{The school year starts in September.}
\end{quote}

A very French notion: after the summer, everyone goes back to school and to work.

\begin{cousinweb}
\emph{rentrée} is built on \emph{rentrer} \textquotedblleft{}to go back in\textquotedblright{}, from \emph{re-} and \emph{entrer}, which is inherited from Latin \emph{intrāre} \textquotedblleft{}to go in\textquotedblright{} (Lewis \& Short).
\end{cousinweb}

\begin{grammarlens}[title={What you've built}]
One more word for telling, explaining and describing at B1.
\end{grammarlens}

\subsection*{Your turn}
\begin{itemize}
\raggedright
  \item \emph{Say it:} \emph{la rentrée}
  \item \emph{Say it:} \emph{la rentrée}, once more
  \item \emph{Say it:} a sentence of your own with \emph{la rentrée}
  \item \emph{Recall:} say the French for the teenager, then the French for youth, then say \emph{la rentrée} again
\end{itemize}

\begin{tcolorbox}[breakable,colback=teal!4,colframe=teal!35!black,title={Before you move on}]
What does la rentrée mean? (\textquotedblleft{}The start of the school year; the return\textquotedblright{}.) Say it once more.
\end{tcolorbox}
\section[majeur]{majeur — of age, adult; major}

\label{lesson:FR-C310-majeur}



\begin{quote}
Before the new one: say the French for youth, then the French for the start of the school year.
\end{quote}

\subsection*{You'll want to know: majeur}
\textbf{majeur} — \textquotedblleft{}of age, adult; major\textquotedblright{}.

\begin{quote}
Il est majeur depuis hier: il a dix-huit ans.
\end{quote}

\begin{quote}
\emph{He came of age yesterday: he is eighteen.}
\end{quote}

In France you are \emph{majeur} at eighteen. \emph{La majeure partie} — most of.

\begin{cousinweb}
\emph{majeur} is a learned borrowing from Latin \emph{māior} \textquotedblleft{}greater\textquotedblright{} (EtymDB), the comparative of \emph{magnus} \textquotedblleft{}great\textquotedblright{}.
\end{cousinweb}

\begin{grammarlens}[title={What you've built}]
One more word for telling, explaining and describing at B1.
\end{grammarlens}

\subsection*{Your turn}
\begin{itemize}
\raggedright
  \item \emph{Say it:} \emph{majeur}
  \item \emph{Say it:} \emph{majeur}, once more
  \item \emph{Say it:} a sentence of your own with \emph{majeur}
  \item \emph{Recall:} say the French for youth, then the French for the start of the school year, then say \emph{majeur} again
\end{itemize}

\begin{tcolorbox}[breakable,colback=teal!4,colframe=teal!35!black,title={Before you move on}]
What does majeur mean? (\textquotedblleft{}Of age, adult; major\textquotedblright{}.) Say it once more.
\end{tcolorbox}
